\chapter{Training Improvements}
\label{ch:improvements}

The experiments reported in Chapter~\ref{ch:experiments} were not the first results produced
by this system.
An iterative development process preceded the final numbers, in which the model's errors were
diagnosed and two targeted interventions were designed and evaluated.
This chapter describes that process: what failures were observed, what interventions were
designed, and how performance changed as a result.

\section{Motivation: Identifying Failure Modes}

After training the initial joint model on the synthetic dataset without augmentation,
three classes of failure were identified through the semantic understanding tests
and per-CWE validation error analysis.

\textbf{Failure 1 --- Safe strong-cryptography functions classified as vulnerable.}
The model systematically classified \texttt{hashlib.sha256} and \texttt{hashlib.sha512}
as vulnerable despite these being the recommended replacements for the genuinely vulnerable
\texttt{hashlib.md5}.
Running the minimal-fix test before any augmentation showed that replacing \texttt{md5}
with \texttt{sha256} left the predicted probability nearly unchanged at 0.906.
The model had learned to associate the presence of \texttt{hashlib} with vulnerability
rather than discriminating between specific algorithm names.

\textbf{Failure 2 --- Safe C functions with \texttt{system()} on hardcoded arguments
classified as vulnerable.}
The minimal-fix test for CWE-78 C (replacing \texttt{system(cmd)} where \texttt{cmd}
contains user input with \texttt{system("ls /safe/path")} where the argument is fully
hardcoded) produced a predicted probability of 0.929 even after the fix.
The training set contained no examples of safe \texttt{system()} usage, so the model
learned \texttt{system(} as an unconditional signal of vulnerability.

\textbf{Failure 3 --- Large synthetic-to-real generalisation gap.}
Experiment~3 on the pre-augmentation model produced a real-world F1 of 0.410 and
a gap of $-0.522$.
Analysis of the missed real-world examples showed that they were structurally more complex
than the synthetic training examples --- longer functions, interleaved control flow,
library-specific idioms --- which the model trained only on short, clean examples could
not reliably recognise.

\section{Intervention 1 --- Hard Negative Mining}

\subsection{Design}

Hard negatives are safe examples deliberately designed to resemble vulnerable ones in
surface form while being semantically safe.
The objective is to force the model to learn the \emph{discriminating} feature rather
than the \emph{correlated} feature.

For the cryptography failure, six hard negative examples were added to the Python safe set:
functions using \texttt{hashlib.sha256}, \texttt{hashlib.sha512}, and
\texttt{hashlib.pbkdf2\_hmac} in the same structural positions where vulnerable examples
use \texttt{hashlib.md5}.
These safe examples use identical code structure to the vulnerable ones, differing only in
the algorithm name --- a minimal contrastive pair that forces the model to attend to the
specific algorithm token and its security implication.

\begin{figure}[H]
\begin{minipage}[t]{0.48\textwidth}
\begin{lstlisting}[language=Python, caption={Vulnerable (CWE-327): \texttt{md5} used as password hash.}]
import hashlib
def hash_password(password):
    return hashlib.md5(
        password.encode()
    ).hexdigest()
\end{lstlisting}
\end{minipage}
\hfill
\begin{minipage}[t]{0.48\textwidth}
\begin{lstlisting}[language=Python, caption={Hard negative (safe): identical structure, secure algorithm.}]
import hashlib
def hash_password(password):
    return hashlib.sha256(
        password.encode()
    ).hexdigest()
\end{lstlisting}
\end{minipage}
\caption{Python CWE-327 hard negative pair: the only difference is \texttt{md5}
vs \texttt{sha256}, forcing the model to discriminate on algorithm name.}
\label{fig:hn_python}
\end{figure}

For the command injection failure, four safe examples were added to the C safe set:
functions calling \texttt{system()} with fully hardcoded string literals as arguments
and no user-controlled data anywhere in the call chain.
Additional safe examples cover OpenSSL SHA-256 and SHA-512 usage (\texttt{SHA256\_Update},
\texttt{SHA512\_Update}).
In total, seventeen new safe examples were added across both languages.

\begin{figure}[H]
\begin{minipage}[t]{0.48\textwidth}
\begin{lstlisting}[language=C, caption={Vulnerable (CWE-78): user input flows into \texttt{system()}.}]
void run_command(char *user_input) {
    char cmd[256];
    sprintf(cmd, "ls %s", user_input);
    system(cmd);
}
\end{lstlisting}
\end{minipage}
\hfill
\begin{minipage}[t]{0.48\textwidth}
\begin{lstlisting}[language=C, caption={Hard negative (safe): \texttt{system()} called with fully hardcoded argument.}]
void run_cleanup(void) {
    system("rm -f /tmp/cache.lock");
}
\end{lstlisting}
\end{minipage}
\caption{C CWE-78 hard negative pair: both call \texttt{system()} but only the
left version has user-controlled data in the call chain.}
\label{fig:hn_c}
\end{figure}

\subsection{Effect}

Table~\ref{tab:hard_neg_effect} shows predicted probabilities for the minimal-fix test
cases before and after adding hard negatives.

\begin{table}[H]
  \centering
  \caption{Minimal-fix test predictions before and after hard negative mining.}
  \label{tab:hard_neg_effect}
  \begin{tabular}{llccc}
    \toprule
    CWE & Version & Before & After & Correct after? \\
    \midrule
    CWE-327 & Vulnerable (md5)       & 0.960 & 0.960 & \checkmark \\
    CWE-327 & Fixed (sha256)         & 0.906 & \textbf{0.048} & \checkmark \\
    CWE-798 & Vulnerable (hardcoded) & 0.960 & 0.960 & \checkmark \\
    CWE-798 & Fixed (env var)        & 0.050 & 0.057 & \checkmark \\
    CWE-89  & Vulnerable (concat)    & 0.959 & 0.959 & \checkmark \\
    CWE-89  & Fixed (parameterised)  & 0.048 & 0.048 & \checkmark \\
    CWE-78  & Vulnerable (user input)& 0.958 & 0.959 & \checkmark \\
    CWE-78  & Fixed (hardcoded path) & 0.929 & 0.929 & $\times$ \\
    \bottomrule
  \end{tabular}
\end{table}

The hard negative examples resolved the CWE-327 failure completely: the probability for a
\texttt{sha256} function dropped from 0.906 to 0.048, flipping the prediction from
VULNERABLE to SAFE.
The CWE-78 C case remained unresolved at 0.929.
This is a fundamentally different failure: there is no surface token in the hardcoded-path
version that distinguishes it from the user-input version --- both call
\texttt{system(cmd)} where \texttt{cmd} is a character buffer populated by \texttt{sprintf}.
A hard negative for this case would require the model to reason about whether the buffer's
value originated from user input, which is a data-flow question that cannot be resolved from
the local function text alone.
This failure is attributed to an architectural limitation rather than a data gap.

Hard negative additions did not degrade overall synthetic F1, which remained at 0.931.

\section{Intervention 2 --- Real-World Data Augmentation}

\subsection{Design}

To reduce the synthetic-to-real generalisation gap, 1{,}000 function-level samples from
the DiverseVul training split (\texttt{sample\_train.jsonl}) were incorporated into the
joint model's training data.
Each real-world record is normalised to the training schema before merging.

The 1{,}000 real samples are concatenated with the full synthetic training set and the
combined data is shuffled before training.
The DiverseVul test split (\texttt{sample\_test.jsonl}, 200 samples) is kept completely
separate; no sample from the test split was exposed to training at any point.

The augmented training set thus contains synthetic examples, hard negatives, and
real-world examples from a different distribution --- a curriculum that exposes the model
to idealised patterns and realistic code within a single fine-tuning pass.

\subsection{Effect on Real-World Performance}

Table~\ref{tab:augmentation_effect} summarises the Experiment~3 results at each stage.

\begin{table}[H]
  \centering
  \caption{Real-world generalisation at each training stage.}
  \label{tab:augmentation_effect}
  \begin{tabular}{lccc}
    \toprule
    Model & Synthetic F1 & Real-world F1 & Gap \\
    \midrule
    Synthetic-only (baseline) & 0.931 & 0.410 & $-0.522$ \\
    + Hard negatives          & 0.931 & 0.471 & $-0.460$ \\
    + Real data augmentation  & 0.931 & \textbf{0.749} & $\mathbf{-0.182}$ \\
    \bottomrule
  \end{tabular}
\end{table}

Hard negatives alone improved real-world F1 from 0.410 to 0.471 (gain: $+0.061$),
primarily by reducing false positives on safe C functions sharing surface tokens with
dangerous ones.
The majority of the improvement came from real-world data augmentation: adding 1{,}000
DiverseVul samples increased real-world F1 by 0.278 points, reducing the generalisation
gap from $-0.460$ to $-0.182$.

Neither intervention degraded synthetic performance; F1 on the held-out synthetic
validation set remained at 0.931 after both interventions.

\subsection{Effect on Error Analysis}

\begin{table}[H]
  \centering
  \caption{Real-world confusion matrix before and after augmentation.}
  \label{tab:cm_before_after}
  \begin{tabular}{lcccccc}
    \toprule
    & TP & FP & TN & FN & FPR & FNR \\
    \midrule
    Pre-augmentation  & 66 & 52 & 32 & 50 & 61.9\% & 43.1\% \\
    Post-augmentation & \textbf{88} & 31 & 53 & \textbf{28} & 36.9\% & 24.1\% \\
    \bottomrule
  \end{tabular}
\end{table}

Both error rates improve substantially.
The false negative rate falls from 43.1\% to 24.1\%: the model now catches substantially
more real-world vulnerabilities.
The false positive rate falls from 61.9\% to 36.9\%: the model is less likely to alarm
on safe real-world code.
The improvement is approximately symmetric, suggesting that the augmentation provided
signal in both directions: real-world vulnerable functions taught the model what genuine
vulnerabilities look like in production code, while the real-world safe functions taught
the model what safe C code looks like beyond the narrow stylistic profile of the synthetic
safe examples.

\section{Discussion: What Each Intervention Contributes}

The two interventions address structurally different problems.
Hard negative mining resolves \emph{within-distribution confusion} by adding contrastive
pairs that anchor the decision boundary at the discriminating token (algorithm name,
argument source), without requiring architectural changes.
Real-world augmentation addresses \emph{out-of-distribution generalisation} by directly
exposing the model to the target distribution --- no synthetic construction can substitute
for real code styles and idioms.
The residual gap of $-0.182$ reflects what neither intervention can fix: unseen CWE classes
(CWE-787, CWE-400, CWE-703) and the architectural inability to track data flow, which
require more training data or a graph-based architecture respectively.
